\section{Introduction}
\label{sec:intro}

A huge amount of music is only available as a printed score, rather than a
machine-readable encoding. Optical Music Recognition (OMR)---the automatic
transformation of scanned images of a score into a symbolic form such as MIDI or
MusicXML---can make these works playable, transposable, searchable, and
accessible~\cite{calvo2020understanding}. Although commercial OMR engines excel
at clear scanner output, the most ubiquitous capture device today is the camera
on a phone, which suffers from perspective distortion, varied lighting, shadows,
and sensor noise. This gap between the clean documents upon which OMR engines are
trained and tested, and the messy photos users provide, is, I argue, the primary
reason OMR is not practical.

In this paper, I measure this gap directly. I build \emph{antiTranscription}, a
pipeline that takes a hand-held photo of monophonic printed sheet music taken
with a phone and produces a MIDI file. It includes both geometric vision
components (rectification, Hough staff detection, morphological segmentation,
geometric pitch inference) and a trained convolutional classifier; this allows
me to distinguish those pipeline components which make a smooth transition from
clean to real photographic images from those that do not.

My experiments highlight a stark dichotomy. While my geometric vision stages
perform robustly (all five photos correctly rectified with sub-degree slant, and
every staff detected), my appearance-based classifier decreases from a top
accuracy of $58.9\%$ on clean PrIMuS glyphs~\cite{calvo2018primus} to $22.7\%$ on
crops from real photographs, a drop of $36.2$ points, which is most pronounced on
thin or hollow glyphs. I show that a segmentation-free pitch reader using only
the geometry of noteheads takes advantage of this by achieving $86\%$ accuracy on
a representative piece, compared to the $14\%$ accuracy of the purely
appearance-based approach.

\paragraph{Contributions.}
(i) A complete five-stage phone-photo-to-MIDI pipeline combining classical
geometry with a learned classifier; (ii) a robust rectification cascade and
Hough/clustering staff detector succeeding on $100\%$ of my real test photos;
(iii) a quantified \emph{clean-to-real domain-gap study} for OMR symbol
classification ($58.9\%\!\rightarrow\!22.7\%$) with per-class analysis of where
appearance models break; and (iv) a geometry-grounded, segmentation-free pitch
reader that improves exact pitch accuracy $6\times$ ($14\%\!\rightarrow\!86\%$)
over that baseline. I scope the system to monophonic treble-clef scores without
accidentals or key signatures, isolating perception from full musical parsing.
